\begin{solapartat}
Si el motor realitza $1,91\cdot 10^{3}$ rpm, la seva freqüència angular serà:
\[ \omega = 1,91\cdot 10^{3}\,\frac{\text{voltes}}{\text{minut}}\times\frac{2\pi\ \text{rad}}{\text{volta}}\times\frac{\text{minut}}{60\text{s}} = 200\ \mathrm{rad/s} \quad\punts{0,2} \]
L'equació del MVHS és:
\[ x(t) = A\,\sin(\omega t + \phi) \]
D'altra banda: $2A = 20,0\ \mathrm{cm}$, per tant $A = 10,0\ \mathrm{cm}$ \quad\punts{0,2}

Com que: $x(t = 0) = \pm A \Rightarrow \phi = \pm\pi/2$ el que és el mateix:
\[ x(t) = A\,\sin(\omega t \pm \pi/2)\ \mathrm{m} \quad\punts{0,2} \]
(també seria vàlida la solució: $x(t) = \pm A\,\cos(\omega t + \phi)$ amb $\phi = 0$)

D'altra banda:
\[ v(t) = \pm A\omega\,\cos(\omega t \pm \pi/2)\ \mathrm{m/s} \quad\punts{0,2} \]
Per tant:
\[ v_\text{màxima} = A\omega = 20,0\ \mathrm{m/s} \quad\punts{0,2} \]
\end{solapartat}

\begin{solapartat}
La constant de recuperació en un MVHS ve donada per: $k = \omega^2 m$ \quad\punts{0,4} de manera que la força recuperadora màxima en un MVHS és:
\[ F_\textit{màxima} = k\,A \quad\punts{0,2} = 800\ \mathrm{N} \quad\punts{0,4} \]
\end{solapartat}
